\chapter{C1 with both $L$ and $\eta$ units}
\label{component:p686209b87229}
\noindent\textbf{Component }\path{p686209b87229}\par
{\small\noindent Source: \path{R13/advances/c1_lunit_eta_unit/eta_unit.tex}\par}
{\small\noindent SHA256: \path{686209b8722928e9065b28ff8b0caeb4dcb38c543c4ff8ae1dc50ddf91744b07}\par}
{\small\noindent Inventory use: LU1--LU12 and cyclic remote join | C14\_\allowbreak{}ENTIRE\_\allowbreak{}ORIGINAL\_\allowbreak{}PARENT\_\allowbreak{}GRAPH; precise subordinate uses retained below\par}
\medskip
\section[C1 with both L and eta units]{C1 with both $L$ and $\eta$ units}
\label{p686209b87229:sec:r7-c1-L-eta-unit}

\begin{theorem}
For the original C1 source with $x\in\Inf_{>0}$ and both
$L=x\zeta$ and $\eta$ positive bounded units, the entire original
parent is tensor-SOS for every independent projective testing pair.
All original positive openings and literal $\beta=0$ are included.
\end{theorem}

The exact source formulas force $D$ strong on this subdomain, so
$D=0$ is absent by source implication.  No opening is removed by
an extra hypothesis, and the proof does not invoke a table or atlas.

\paragraph{The source and the only testing strip left by four events.}
Use the exact shared geometry (LC1), and write $\rho=r_n$ for the
native shape parameter and $r=1-\theta$ for the geometric one.
Thus
\[
 B=L,\quad y=B\eta,\quad H=B\rho=B-yt,\quad
 t=(1-\rho)/\eta,\quad j=1+zr-t,
 \quad z=x\in\Inf_{>0}.
\]
Here $B,y,\eta$ are units.  For the boundary coefficient formulas
below set
\[
 P=1+B+y,\qquad P_4=1+B(1-\rho),\qquad S=H+y=P-P_4.
                                                               \tag{LU1}
\]
All three quantities are positive units, uniformly through
$\rho\to0$, $1-\rho\to0$, $j\to0$, and $1-H\to0$.
The native formulas imply that $D,C,A,\alpha$ are strong.
One forward cyclic relabeling sends these four consecutive events
to $C,A,\alpha,m$ and sends testing coordinates to
$(u',v')=(v,1/u)$.  The full-parent four-consecutive theorem
(CE11)--(CE16) therefore covers every test except
\[
                    u\in\Inf,\qquad v\text{ a bounded nonzero unit}.
                                                               \tag{LU2}
\]
The cycle is the inverse of (CE9)--(CE10), so its literal endpoint
scope is the original homogeneous parent's scope.

It remains to prove (LU2).  All roots and coefficients used below
are valid for both testing signs, including literal $u=0$.

\paragraph{A bounded full-parent remainder with every opening retained.}
Use the graph, harmonic roots and gauge (LII2)--(LII8), with
$k=\sqrt y$ now a source unit.  The bounded source data are
\[
 c=Cyt,\quad d_p=Dyj,\quad
 \mu=m/z^2,\quad \sigma=\alpha(1-H),\quad
 \tau=\beta(1-r),\quad A_H=AH.
\]
They are bounded; $c,d_p,\mu,\sigma,A_H$ are units.  In residue,
with the boundary symbols of (LU1), their exact values are
\[
 c=\frac S{2P},\quad d_p=\frac{BP_4}{2P},\quad
 \mu=\frac{Sr(1+y)}{2B^2P},\quad
 \sigma=\frac{BS}{2P},\quad A_H=\frac{yP_4}{2P},\quad
 \tau=\frac S{2P}-\frac{1-r}{2}.                         \tag{LU3}
\]
These are residue identities, not substitutions of source and
testing fields.  They are obtained after multiplying away every
possibly small opening in the exact native formulas.
Since $\tau\ge0$, $r\ge P_4/P$ in residue and is a positive unit.
Also $t+j\sim1$, so $\ell=y(C+D)$ has a positive lower bound,
$q=D/(C+D)$ is bounded in $[0,1]$, and
$g=yCD/(C+D)$ is a positive unit.

Choose positive bounded caps $\ell_*\le\ell$,
$A_*\le A$, and $\alpha_*\le\alpha$.  Choose a bounded cap
$0\le\beta_*\le\beta$ as specified below.  Retain the five
contact-centered positive packets
\[
\begin{gathered}
 4(\ell-\ell_*)N^2,\qquad 4CV\eta_t^2,\qquad
 4(A-A_*)D_uD_v\eta_t^2,\qquad
 4(\alpha-\alpha_*)D_u(v\eta_t-\zeta)^2,\\
 4(\beta-\beta_*)D_v(\xi+u\zeta+R/B)^2,
 \qquad \eta_t=k(N+qR_h),\quad R=v(U-z)W.
\end{gathered}
\tag{LU4}
\]
The exact remainder is (LII6), with its $n$ line deleted, its $A$
packet replaced by
\[
 4A_*D_uD_v\eta_t^2
 +\frac{8A_H}B D_uD_vuW\eta_t
 +\frac{4A_HH}{B^2}D_uD_vUW^2,
                                                               \tag{LU5}
\]
and with (LC4) plus $4\beta_*D_v(\xi+u\zeta+R/B)^2$ added.
This specifies every coefficient of the full remainder.  Only
bounded source products occur, even at an unbounded original weight.

The same dependent-root gauge and exact contact equations give
(LII9)--(LII13).  The new constant scalar coefficients are
\[
 \frac{d_s+e_sV}{z^2}=
 \frac4{B^2}\left[
 \mu(B+zr)^2+
 \left\{\frac{r^2}{2}-\tau(1+r)+\beta_*\right\}(1+V)
 \right].                                                   \tag{LU6}
\]
The support $P_\xi\in(U,z)$ and scalar degree at most two in $U$
remain exact.  Apply all of (LII15).  As $v$ is a unit, eliminate
$T$ by the exact pure-factor identity
\[
                        T=(uk/v)R.                         \tag{LU7}
\]
Its coefficient is a bounded infinitesimal tensor.  Thus the entire
remainder has a bounded Gram on
$(N,\zeta,R_h,\xi,R)$.  No source/testing ratio or mixed inverse
has been introduced.  We now verify its complete residue.

\paragraph{A three-square identity supplies all three state/angle normals.}
The following boundary computation is an exact rational identity
before taking any opening limit.  Define
\[
\begin{aligned}
 C_0&=\frac{\rho+\eta}{2(1-\rho)P},&
 A_0&=\frac{\eta P_4}{2\rho P},\\
 D_0&=\frac{P_4}{2(\eta-1+\rho)P},&
 \alpha_0&=\frac{B^2(\rho+\eta)}{2(1-B\rho)P}.
\end{aligned}
\]
These may have unbounded openings; only the final capped identity
below is used as a bounded Gram.  For an indeterminate $p$, put
\[
 e=Y-\rho p,\qquad f=Z-vp/B,\qquad T_v=(V-y)p.
\]
The entire leading three-state form is exactly
\[
\begin{split}
 &4C_0D_v(p-Y)^2+4A_0D_vY^2
 +(4D_0+2)(p-Y-vZ)^2+(4\alpha_0+2)(vY-Z)^2-2D_vZ^2\\
 &\quad=\frac{4D_0}{B^2}(T_v+Be+Bvf)^2
       +4\alpha_0(f-ve)^2+4(C_0+A_0+\tfrac12)D_ve^2.
\end{split}                                                    \tag{LU8}
\]
This identity holds for every $v$, not only near the actual Bob
angle.  In the contact allocation at $z=u=0$,
$e=k(N+qR_h)$ and $T_v=-B(kR_h+v\zeta)$.
Subtract the residue contributions of the retained packets (LU4).
The harmonic identity then gives the complete three-normal block
\[
 \boxed{4\ell_*N^2+4gR_h^2
 +4(A_*+\tfrac12)D_vy(N+qR_h)^2
 +4\alpha_*\{\zeta-vk(N+qR_h)\}^2.}                     \tag{LU9}
\]
It has strictly positive residue on $(N,\zeta,R_h)$ because
$\ell_*,g,\alpha_*$ are positive units.  No normal has been
suppressed.  Every rational cancellation used to obtain (LU9)
has bounded source products on both sides, so the identity extends
continuously through every opening listed after (LU1).

\paragraph{The remaining two-root block and a finite cap.}
There is no finite residue coupling between these three roots and
$(\xi,R)$.  Such terms carry $u$, $z$, or the coefficient in (LU7).
For completeness, the exact first-normal support gives this block
without dropping a term.  To first order in $(U,z)$ write
$P_\xi=UA(V)+zZ(V)$; both coefficient functions are affine in $V$,
and
\[
 Z(V)=\frac{D_v}{B}(-4r+8\tau-8\beta_*).
\]
The exact contact equation gives $A(y)+Z(y)=0$ in residue.
After the first allocation of (LII15), the coefficient of $\xi R$
is therefore $A(V)-(V-y)(A'+Z')=-Z(V)$.
Using (LU3) and (LU6), the complete block is
\[
 M_{\xi R}=
 \begin{pmatrix}
 2D_v(1+2\beta_*)&\dfrac{2D_v}{B}(P_4/P+2\beta_*)\\[1mm]
 \dfrac{2D_v}{B}(P_4/P+2\beta_*)&
 \dfrac2{B^2}\{(P_4/P+2\beta_*)D_v-(S/P)r(V-y)\}
 \end{pmatrix}.                                           \tag{LU10}
\]
Here all letters denote real residues when positivity is assessed.

If $\beta$ is bounded, take $\beta_*=\beta$.  The native relation
$(1-r)(2\beta+1)=S/P$ follows from (LU3), and (LU10) has determinant
\[
 \frac{4S D_v}{B^2P}(1+2\beta)r(1+y)>0.                 \tag{LU11}
\]
Its first diagonal entry is positive, including $\beta=0$.
If $\beta$ is unbounded, take $\beta_*$ to be a sufficiently large
positive real cap.  Increasing it adds the positive square
$4\beta_*D_v(\xi+R/B)^2$.  On its kernel $\xi=-R/B$,
the uncapped block equals
\[
          \frac{2S}{B^2P}\{1+ry+(1-r)V\}R^2>0.           \tag{LU12}
\]
The elementary real two-by-two Schur argument therefore supplies
a finite cap making (LU10) strictly positive.  Unboundedness of the
original $\beta$ permits that cap.

Together (LU9)--(LU12) prove strict positive residue of the full
five-root Gram.  The bounded Gram lemma gives its tensor SOS;
restoring (LU4) proves (LU2).  The cyclic four-event complement
completes every testing chart and original projective endpoint.

\paragraph{Verification scope.}
The checker verifies the exact three-square state identity, all three
capped normals, both explicit two-root positivity formulas, the scalar
coefficient and support assertions, and prints the full extracted
first-normal and scalar polynomials.  The exact source contact and
monic contact divisions are the unchanged full-parent argument of
(LII7)--(LII15), with the bounded modifications (LU4)--(LU6).
This theorem closes the $L$-unit, $\eta$-unit subdomain of C1.4;
the $\eta$-unbounded subdomain is not included.
